\documentclass[12pt]{article}
\usepackage[utf8]{inputenc}
\usepackage[T1]{fontenc}
\usepackage[margin=1in]{geometry}
\usepackage{setspace}
\usepackage{xcolor}
\usepackage[colorlinks=true, urlcolor=blue!60!black, linkcolor=blue!60!black]{hyperref}
\setlength{\parindent}{0em}
\setlength{\parskip}{0.6em}
\renewcommand{\baselinestretch}{1.0}
\begin{document}
\begin{center}
{\LARGE \textit{The Civil Rights Cases}, 109 U.S. 3 (1883)\par}
\vspace{0.5em}
{\large Case Brief --- Two-Page Edition\par}
\end{center}

\section{Memory Jogger}
The Fourteenth Amendment reaches only state action, so Congress could not ban private racial discrimination in inns, theaters, and railroads under the Civil Rights Act of 1875 (\href{https://www.courtlistener.com/opinion/90897/civil-rights-cases/}{CourtListener}).

\section{Detailed Case Facts}
The 1875 Act guaranteed all persons ``full and equal enjoyment'' of inns, public conveyances, and theaters and punished any person who violated it. Five consolidated cases involved innkeepers, theaters in San Francisco and New York, and a railroad conductor who excluded Sallie Robinson, a Black woman, from the ladies' car.

\section{Procedural History}
Three circuit courts certified divided opinions; in two others, the losing parties sought review by writ of error.

\section{Judicial Votes}
8--1. Justice Bradley wrote for the Court; Justice Harlan dissented.

\section{Holding}
Sections 1 and 2 are void. Section 5 of the Fourteenth Amendment permits only corrective legislation against state action, and private denial of accommodations is not a badge or incident of slavery under the Thirteenth Amendment. The Commerce Clause question was reserved.

\section{Analysis of Opinions}
\textbf{Majority:} The Act regulated private individuals directly, even in States with fair laws, creating a federal ``code of municipal law''; calling every private refusal a badge of slavery would be ``running the slavery argument into the ground.'' \textbf{Harlan:} the reasoning was ``entirely too narrow and artificial''; railroads, inns, and licensed theaters serve public functions as ``agents or instrumentalities of the State,'' and Congress once had ample power to protect slaveholders.

\section{Examples: Future Applications}
\textbf{Same-Side:} a federal remedy against private gender-motivated violence fails (\textit{United States v. Morrison}). \textbf{Opposite-Side:} court enforcement of racial covenants is state action (\textit{Shelley v. Kraemer}). \textbf{Fence-Sitter:} a segregating restaurant leased in a city-owned garage.

\section{Critique}
Defenders cite the text and federalism; originalists such as Michael McConnell and Randy Barnett argue Reconstruction intended more. Progressives place the case in the anti-canon, and Charles Black called the public/private line illusory.

\section{Key Quotations}
\begin{quote}
``It is State action of a particular character that is prohibited. Individual invasion of individual rights is not the subject matter of the amendment.''
\end{quote}
\noindent --- Justice Bradley, Opinion of the Court, 109 U.S. at 11.

\end{document}
